% Source: https://artofproblemsolving.com/community/c3h1884155_subset_i_just_have_a_couple_questions_you_dont_have_to_solve_a_problem
% Stable identity: aops-topic-1884155-q01
(a) Yes. A set is determined only by which elements it contains; the order of listing (and repetition) does not matter, so the two descriptions are the same set. In particular this set counts once among its own subsets.

(b) First take subsets of $\{-2,-1,1,2\}$ with sum $0$. Such a subset consists of a positive part $P\subseteq\{1,2\}$ and a negative part $N\subseteq\{-1,-2\}$ with opposite sums. The possible sums of $P$ are $0,1,2,3$ (for $\varnothing,\{1\},\{2\},\{1,2\}$), each from exactly one $P$, and the sums of $N$ are $0,-1,-2,-3$, each from exactly one $N$. So there are exactly four zero-sum subsets:
\[
\varnothing,\quad\{-1,1\},\quad\{-2,2\},\quad\{-2,-1,1,2\}.
\]
Each can be taken with or without $0$, giving $8$ subsets with sum $0$; removing the empty set leaves $\boxed{7}$:
\[
\{0\},\ \{-1,1\},\ \{-2,2\},\ \{-1,0,1\},\ \{-2,0,2\},\ \{-2,-1,1,2\},\ \{-2,-1,0,1,2\}.
\]
